\documentclass[11pt,a4paper]{article}
\usepackage[T1]{fontenc}
\usepackage{lmodern}
\usepackage[margin=27mm]{geometry}
\usepackage{amsmath,amssymb,amsthm,mathtools}
\usepackage{microtype}
\usepackage{booktabs,array}
\usepackage{enumitem}
\usepackage[hidelinks]{hyperref}
\setlength{\parskip}{5pt plus 1pt}
\setlength{\parindent}{0pt}
\setlength{\emergencystretch}{3em}
\setlist{itemsep=2pt,topsep=3pt}
\raggedbottom
\newtheorem{theorem}{Theorem}[section]
\newtheorem{lemma}[theorem]{Lemma}
\newtheorem{corollary}[theorem]{Corollary}
\newtheorem{proposition}[theorem]{Proposition}
\theoremstyle{definition}
\newtheorem{definition}[theorem]{Definition}
\newtheorem{example}[theorem]{Example}
\theoremstyle{remark}
\newtheorem{remark}[theorem]{Remark}
\newcommand{\Z}{\mathbb{Z}}
\newcommand{\N}{\mathbb{N}}
\newcommand{\Q}{\mathbb{Q}}
\newcommand{\F}{\mathbb{F}}
\newcommand{\ps}[1]{#1[[x]]}
\newcommand{\coeff}[2]{[x^{#1}]\,#2}
\newcommand{\hl}[1]{\textbf{#1}}
\newcommand{\para}[1]{\par\medskip\noindent\textbf{#1}\enspace}
\hypersetup{pdftitle={The odd-indexed terms of a nested square are even after the first term: a proof of an observation in OEIS A392210},pdfauthor={Dev D. Goyal},pdfsubject={OEIS A392210}}
\title{The odd-indexed terms of a nested square are even after the first term\\[6pt]\large A proof of an observation in OEIS A392210}
\author{Dev D. Goyal}
\date{October 2026}
\begin{document}
\maketitle
\begin{abstract}
The sequence A392210 of the On-Line Encyclopedia of Integer Sequences (OEIS), contributed by Paul D.~Hanna, is defined
by an ``infinitely nested square''
\[
x+\mathrm{sq}\Bigl(\frac{x}{1-x}+\mathrm{sq}\Bigl(\frac{x}{1-2x}+\mathrm{sq}\Bigl(\frac{x}{1-3x}+\cdots\Bigr)\Bigr)\Bigr),
\qquad \mathrm{sq}(y)=y^2 .
\]
Its terms are $1,1,4,18,96,579,\dots$, and the entry observes that the terms $a(2n-1)$ with $n>1$ appear to be even.
We first make the nested expression precise: cutting it off after $d$ levels gives the correct coefficients up to
degree $d+1$, and the resulting series is the unique solution of the simpler equation $A(x)=x+A\bigl(x/(1-x)\bigr)^2$.
We then prove the observation: once all coefficients are reduced modulo $2$, the square of any power
series contains only even powers of $x$. The same idea also gives a recursive rule for the parity of the even-indexed
terms. The main results have been checked by computer and proved formally in the Lean proof assistant.
\end{abstract}
\begin{center}\begin{minipage}{0.88\linewidth}\small
\textbf{How this paper was produced.} This paper reports the result of an AI research effort directed by
Dev~D.~Goyal. ChatGPT (OpenAI), using the GPT-6 model with Pro reasoning and working autonomously from prompts by
Dev~D.~Goyal, found the proof, wrote its Lean~4 formalization, and wrote a first draft of the paper.
Dev~D.~Goyal chose the problems and asked for proofs, formal verification and revisions, but did not supply proof
ideas. Claude Opus~5.5, used with Claude Code, then checked the work independently: it rebuilt the formal proofs from
source and confirmed that they rely only on Lean's three standard axioms, recomputed the sequence from its definition
compared it with the published terms and tested every claim numerically, and read every step of the proofs. It also wrote the present
version of this paper.
Dev~D.~Goyal has verified the results and is responsible for them. The paper has not been peer reviewed.
The section ``About this paper'' at the end gives the details.
\end{minipage}\end{center}

\section{Introduction}

\subsection{The sequence}\label{sec:sequence}
The \emph{On-Line Encyclopedia of Integer Sequences} (OEIS)~\cite{oeisdb} is a large public database of integer
sequences. Many
entries also record observed patterns that the contributor did not prove.

This paper is about entry A392210~\cite{oeis}, contributed by Paul~D.~Hanna in February 2026. The entry reads, in part:
\begin{quote}\small\ttfamily\raggedright G.f.: x + sq(x/(1-x) + sq(x/(1-2*x) + sq(x/(1-3*x) + sq(x/(1-4*x) + ...)))) where sq(x) = x\^{}2, an infinitely nested square.\par\smallskip This sequence equals the self-convolution (shifted right) of the binomial transform of this sequence.\par\smallskip It appears that a(2*n-1) is even for n > 1.\par\smallskip (1) A(x) = x + A(x/(1-x))\^{}2.\par\smallskip (5) A(x/(1-n*x)) = x/(1-n*x) + A( x/(1-(n+1)*x) )\^{}2 for all n.\par\smallskip\hfill\textrm{--- OEIS A392210, Paul D. Hanna, February 2026}\end{quote}

Here ``G.f.''\ stands for \emph{generating function}: the terms $a(1),a(2),a(3),\dots$ of the sequence are the
coefficients of the power series $A(x)=a(1)\,x+a(2)\,x^2+a(3)\,x^3+\cdots$. The expression $\frac{x}{1-x}$ stands for
the power series $x+x^2+x^3+\cdots$, and more generally $\frac{x}{1-cx}=x+cx^2+c^2x^3+\cdots$.
Section~\ref{sec:background} explains these notions. To avoid confusion with the equation numbers of this paper,
we refer to the formulas numbered (1) and (5) in the entry as (E1) and (E5).

The first line defines the series by an \emph{infinitely nested} expression: a square inside a square inside a square,
without end. Such an expression has to be given a meaning. Section~\ref{sec:results} gives a precise meaning, by
cutting the nesting off after finitely many levels, and shows that the result is the same as the series defined by
formula (E1) of the entry,
\[
A(x)=x+A\Bigl(\frac{x}{1-x}\Bigr)^2 .
\]
Here $A\bigl(\frac{x}{1-x}\bigr)$ is the series obtained by substituting $\frac{x}{1-x}$ for $x$ in $A$, and
$A\bigl(\frac{x}{1-x}\bigr)^2$ is its square. We take formula (E1) as the definition.

The second line of the quotation says the same thing as formula (E1), in words. Substituting $\frac{x}{1-x}$ into $A$
gives the series whose coefficients are
\[
b(n)=\sum_{k=1}^{n}\binom{n-1}{k-1}\,a(k),
\]
which is called the \emph{binomial transform} of the sequence. Squaring a series combines its coefficients into the sums
$\sum_kb(k)\,b(n-k)$, called the \emph{self-convolution}. So formula (E1) says that $a(1)=1$ and, for $n>1$,
\[
a(n)=\sum_{k=1}^{n-1}b(k)\,b(n-k),
\]
which lets one compute the terms one at a time. The first terms, starting at $n=1$, are
1, 1, 4, 18, 96, 579, 3834, 27404, 209052, 1688389, \dots

\subsection{The question}
The entry observes that $a(2n-1)$ appears to be even for every $n>1$, that is, that $a(3),a(5),a(7),\dots$ are all even.
The table shows the parity of the first terms ($1$ means odd, $0$ means even). The odd indices $n\ge3$ are printed in
bold.
\begin{center}\small\setlength{\tabcolsep}{3.2pt}
\begin{tabular}{>{\raggedright\arraybackslash}p{6.4em}>{\centering\arraybackslash}p{1.15em}>{\centering\arraybackslash}p{1.15em}>{\centering\arraybackslash}p{1.15em}>{\centering\arraybackslash}p{1.15em}>{\centering\arraybackslash}p{1.15em}>{\centering\arraybackslash}p{1.15em}>{\centering\arraybackslash}p{1.15em}>{\centering\arraybackslash}p{1.15em}>{\centering\arraybackslash}p{1.15em}>{\centering\arraybackslash}p{1.15em}>{\centering\arraybackslash}p{1.15em}>{\centering\arraybackslash}p{1.15em}>{\centering\arraybackslash}p{1.15em}>{\centering\arraybackslash}p{1.15em}>{\centering\arraybackslash}p{1.15em}>{\centering\arraybackslash}p{1.15em}}
\toprule
$n$ & 1 & 2 & \hl{3} & 4 & \hl{5} & 6 & \hl{7} & 8 & \hl{9} & 10 & \hl{11} & 12 & \hl{13} & 14 & \hl{15} & 16 \\ \midrule
$a(n)\bmod 2$ & 1 & 1 & \hl{0} & 0 & \hl{0} & 1 & \hl{0} & 0 & \hl{0} & 1 & \hl{0} & 1 & \hl{0} & 1 & \hl{0} & 1 \\
\bottomrule
\end{tabular}
\\[4pt]
\begin{tabular}{>{\raggedright\arraybackslash}p{6.4em}>{\centering\arraybackslash}p{1.15em}>{\centering\arraybackslash}p{1.15em}>{\centering\arraybackslash}p{1.15em}>{\centering\arraybackslash}p{1.15em}>{\centering\arraybackslash}p{1.15em}>{\centering\arraybackslash}p{1.15em}>{\centering\arraybackslash}p{1.15em}>{\centering\arraybackslash}p{1.15em}>{\centering\arraybackslash}p{1.15em}>{\centering\arraybackslash}p{1.15em}>{\centering\arraybackslash}p{1.15em}>{\centering\arraybackslash}p{1.15em}>{\centering\arraybackslash}p{1.15em}>{\centering\arraybackslash}p{1.15em}>{\centering\arraybackslash}p{1.15em}>{\centering\arraybackslash}p{1.15em}}
\toprule
$n$ & \hl{17} & 18 & \hl{19} & 20 & \hl{21} & 22 & \hl{23} & 24 & \hl{25} & 26 & \hl{27} & 28 & \hl{29} & 30 & \hl{31} & 32 \\ \midrule
$a(n)\bmod 2$ & \hl{0} & 1 & \hl{0} & 1 & \hl{0} & 1 & \hl{0} & 0 & \hl{0} & 1 & \hl{0} & 1 & \hl{0} & 1 & \hl{0} & 1 \\
\bottomrule
\end{tabular}
\\[4pt]
\begin{tabular}{>{\raggedright\arraybackslash}p{6.4em}>{\centering\arraybackslash}p{1.15em}>{\centering\arraybackslash}p{1.15em}>{\centering\arraybackslash}p{1.15em}>{\centering\arraybackslash}p{1.15em}>{\centering\arraybackslash}p{1.15em}>{\centering\arraybackslash}p{1.15em}>{\centering\arraybackslash}p{1.15em}>{\centering\arraybackslash}p{1.15em}>{\centering\arraybackslash}p{1.15em}>{\centering\arraybackslash}p{1.15em}>{\centering\arraybackslash}p{1.15em}>{\centering\arraybackslash}p{1.15em}>{\centering\arraybackslash}p{1.15em}>{\centering\arraybackslash}p{1.15em}>{\centering\arraybackslash}p{1.15em}>{\centering\arraybackslash}p{1.15em}}
\toprule
$n$ & \hl{33} & 34 & \hl{35} & 36 & \hl{37} & 38 & \hl{39} & 40 & \hl{41} & 42 & \hl{43} & 44 & \hl{45} & 46 & \hl{47} & 48 \\ \midrule
$a(n)\bmod 2$ & \hl{0} & 1 & \hl{0} & 1 & \hl{0} & 1 & \hl{0} & 0 & \hl{0} & 1 & \hl{0} & 1 & \hl{0} & 1 & \hl{0} & 0 \\
\bottomrule
\end{tabular}
\end{center}

\subsection{What this paper proves}
\textbf{Main results.} \emph{(i) The nested expression in the name of the entry has a precise meaning, and it defines the same
series as formula (E1) (Theorem~\ref{thm:def}). (ii) The coefficient $a(2n-1)$ is even for every $n>1$
(Theorem~\ref{thm:main}).}

The proof of (ii) is short. Modulo $2$, the square of a power series contains only even powers of $x$. So, by formula
(E1), modulo $2$ the only odd power of $x$ in $A(x)$ is the single term $x$.

\subsection{What was known}
The entry lists eight equivalent formulas for $A$, among them formula (E1), formula (E5), and the nested expression in
its name, and states the evenness of $a(2n-1)$ as an observation (``It appears that''). The observation follows quickly
from formula (E1); we are not aware of a written proof. We include it
for completeness, together with a precise treatment of the nested expression.

\subsection{Motivation}
Infinitely nested expressions, such as nested squares, nested square roots, or continued fractions, are a common way of
describing power series in the OEIS. They raise a natural question: what does the expression mean? For power series
there is a general answer, which we explain here: each additional level of nesting fixes one more
coefficient, so the coefficients of the truncated expressions eventually stop changing. The parity argument, in turn, is
a basic tool that is used throughout this collection of papers: modulo $2$, squaring a series doubles all its
exponents.

\subsection{How to read this paper}
Section~\ref{sec:background} explains the notions we need: power series, substitution, the fractions
$\frac{x}{1-cx}$, and remainders. Section~\ref{sec:results} states the results precisely, and Section~\ref{sec:proofs}
proves them. Section~\ref{sec:checks} describes how the results were checked, by computer and by a formal proof in
Lean. Section~\ref{sec:discussion} discusses the scope of the result, and the parity of the even-indexed terms.
\section{Background}\label{sec:background}
This section collects the definitions and facts used later. Readers who know formal power series and congruences can skip it and return when needed.
\subsection{Formal power series}
A \emph{formal power series} is an expression
\[
U(x)=u_0+u_1x+u_2x^2+u_3x^3+\cdots ,
\]
determined by its sequence of coefficients $u_0,u_1,u_2,\dots$; see~\cite{niven,stanley}. No number is ever
substituted for $x$, so questions of convergence do not arise. We write $\coeff{n}{U}=u_n$ for the coefficient of $x^n$ in $U$; two series are equal
when all their coefficients are equal.

Series are added coefficient by coefficient and multiplied like polynomials:
\[
\coeff{n}{(UV)}=u_0v_n+u_1v_{n-1}+\cdots+u_nv_0 .
\]
Each coefficient of a product is a finite sum. For example, $(1-x)(1+x+x^2+x^3+\cdots)=1$, so $1+x+x^2+\cdots$ is
the \emph{multiplicative inverse} of $1-x$, written $\frac{1}{1-x}$. More generally, every series $D$ with constant term
$1$ has a unique multiplicative inverse $\frac1D$, whose coefficients are determined one at a time by $D\cdot\frac1D=1$
using only addition, subtraction and multiplication. Whenever a quotient of power series appears in this paper, its
denominator has constant term $1$, so the quotient is again a power series.

A \emph{commutative ring} is a set with an addition, a subtraction and a multiplication that satisfy the usual rules
of arithmetic; examples are $\Z$, $\Q$ and the integers modulo $m$. For a commutative ring $R$ we write $\ps{R}$ for
the set of power series with coefficients in $R$; for instance, writing $\F_2$ for the integers modulo $2$, $\ps{\F_2}$
consists of the series with coefficients modulo $2$. We write
$U(0)$ for the constant term $u_0$ of $U$, and $U=O(x^d)$ when the coefficients of $1,x,\dots,x^{d-1}$ in $U$ are all
zero. Thus $U-V=O(x^d)$ means that $U$ and $V$ have the same coefficient of $x^n$ for every $n<d$. A series with
constant term $0$ and coefficient of $x$ equal to $1$, that is, a series of the form $x+u_2x^2+u_3x^3+\cdots$, is
called \emph{normalized}.

\subsection{Substituting one series into another}
If $V(x)$ has constant term $0$, we can substitute $V$ for $x$ in $U$:
\[
U(V(x))=u_0+u_1V(x)+u_2V(x)^2+u_3V(x)^3+\cdots .
\]
This is called \emph{composition}, also written $U\circ V$. It makes sense because $V(x)^k$ has no terms below $x^k$,
so only the finitely many terms with $k\le n$ affect the coefficient of $x^n$. For example, if $U=x+x^2$ and $V=x^2$, then
$U(V(x))=x^2+x^4$. Composition is associative: $U\circ(V\circ W)=(U\circ V)\circ W$.

Substituting a series into itself repeatedly gives its \emph{iterates}: $U^{\circ2}(x)=U(U(x))$,
$U^{\circ3}(x)=U(U(U(x)))$, and so on, with $U^{\circ0}(x)=x$ and $U^{\circ1}=U$. They must not be confused with
\emph{powers} $U^2=U\cdot U$, $U^3=U\cdot U\cdot U$. (The OEIS sometimes writes $A^k(x)$ for the $k$-th iterate. In
this paper the small circle always means an iterate, and a plain exponent always means a power.)

\subsection{The fractions \texorpdfstring{$x/(1-cx)$}{x/(1-cx)}}
For a number $c$ we write
\[
M_c(x)=\frac{x}{1-cx}=x+cx^2+c^2x^3+c^3x^4+\cdots .
\]
In particular $M_0(x)=x$. These series satisfy a simple rule under substitution:
\begin{equation}\label{eq:mobius}
M_c\bigl(M_d(x)\bigr)=M_{c+d}(x).
\end{equation}
To see this, start from $M_c(y)\,(1-cy)=y$, which is the definition of $M_c$, and substitute $y=M_d(x)$. Since
$M_d(x)\,(1-dx)=x$, multiplying both sides by $1-dx$ gives $M_c(M_d(x))\,(1-dx-cx)=x$, which says exactly that
$M_c(M_d(x))=\frac{x}{1-(c+d)x}$. (Here we used that substituting into a product is the same as multiplying the
substituted factors.) For example, substituting $\frac{x}{1-x}$ into itself gives $\frac{x}{1-2x}$.

Substituting $M_1$ into a series is the binomial transform mentioned in the introduction: since
$M_1(x)^k=x^k(1-x)^{-k}=\sum_{n\ge k}\binom{n-1}{k-1}x^n$, the coefficient of $x^n$ in $A(M_1(x))$ is
$\sum_k\binom{n-1}{k-1}a(k)$.

\subsection{Remainders and congruences}
For an integer $m\ge2$ we write $a\equiv b\pmod m$ (``$a$ is congruent to $b$ modulo $m$'') when $a-b$ is divisible by
$m$, that is, when $a$ and $b$ leave the same remainder after division by $m$. The remainder of $a$ on division by
$m$ is called its \emph{residue} modulo $m$; it is one of $0,1,\dots,m-1$.

Addition and multiplication are compatible with congruences, so the residues $0,1,\dots,m-1$, added and multiplied
modulo $m$, form a ring, written $\Z/m\Z$. When $m=p$ is a prime we also write $\F_p$. In $\F_2=\{0,1\}$ we have $1+1=0$,
so addition and subtraction coincide.

For power series with integer coefficients, $U\equiv V\pmod m$ means that \emph{every} coefficient of $U-V$ is
divisible by $m$. Equivalently, the two series become equal when every coefficient is replaced by its residue; we call
this \emph{reducing} the series modulo $m$. The key fact is the following.

\emph{Reduction modulo $m$ is compatible with all the operations above.} To reduce a sum, a product, a substitution, or
an inverse, we may reduce the ingredients first and then compute in $\Z/m\Z$. This holds because every
coefficient of the result is obtained from finitely many coefficients of the ingredients by adding, subtracting and
multiplying, and these operations are compatible with remainders.

\para{Rings.} Some statements below hold over every commutative ring, for example $\Z$, $\Q$ and $\Z/m\Z$. No
division is used in them; this matters because in $\Z/m\Z$ division is not always possible (for instance, $2$ has no inverse in
$\Z/4\Z$).

\subsection{The Frobenius identity}
\begin{lemma}[Frobenius identity]\label{lem:frobenius}
Let $p$ be a prime. For every power series $U$ with coefficients in $\F_p$, we have $U(x)^p=U(x^p)$.
\end{lemma}

\begin{proof}
In $\F_p$, the binomial theorem simplifies to $(a+b)^p=a^p+b^p$, because every binomial coefficient
$\binom pi$ with $0<i<p$ is divisible by $p$. Also $c^p=c$ for every $c\in\F_p$ (Fermat's little theorem). So
$\bigl(\sum_nu_nx^n\bigr)^p=\sum_nu_n^px^{pn}=\sum_nu_nx^{pn}=U(x^p)$. The first equality holds for polynomials by
repeated use of $(a+b)^p=a^p+b^p$; it holds for power series because each coefficient of $U(x)^p$ involves only finitely
many coefficients of $U$.
\end{proof}

For $p=2$ this says that squaring a series modulo $2$ doubles every exponent: for instance,
$(x+x^3)^2\equiv x^2+x^6\pmod2$. In particular, the square of any series has, modulo $2$, no terms of odd degree.

\section{Results}\label{sec:results}

We take formula (E1) of the entry as the definition:
\begin{equation}\label{eq:def}
A(x)=x+A\bigl(M_1(x)\bigr)^2,\qquad M_1(x)=\frac{x}{1-x}.
\end{equation}

\para{The nested expression.} To give the nested expression a precise meaning, cut the nesting off after $d$ levels. That is, for
$d\ge1$ define series $H^{(d)}_d,H^{(d)}_{d-1},\dots,H^{(d)}_1$ from the inside out by
\[
H^{(d)}_d=M_d,\qquad H^{(d)}_j=M_j+\bigl(H^{(d)}_{j+1}\bigr)^2\quad(j=d-1,\dots,1),
\]
and put $A^{[d]}=x+\bigl(H^{(d)}_1\bigr)^2$. For example,
\[
A^{[3]}=x+\Bigl(M_1+\bigl(M_2+M_3^2\bigr)^2\Bigr)^2 .
\]
These are power series, obtained by finitely many additions, multiplications and divisions by $1-cx$.

\begin{theorem}[What the nested square means]\label{thm:def}
\begin{enumerate}[label=(\alph*)]
\item Equation~\eqref{eq:def} has exactly one normalized solution $A$ with integer coefficients (and, more generally,
exactly one with coefficients in any commutative ring).
\item For $j\ge0$ put $G_j=A(M_j(x))$. Then $G_0=A$ and
\[
G_j=M_j+G_{j+1}^2\qquad\text{for every }j\ge0 .
\]
This is formula (E5) of the entry.
\item For every $d\ge1$, the series $A^{[d]}$ and $A$ have the same coefficients up to degree $d+1$:
$A-A^{[d]}=O(x^{d+2})$.
\end{enumerate}
\end{theorem}

Part (c) says that the coefficient of $x^n$ in $A^{[d]}$ no longer changes once $d\ge n-1$. The infinitely nested
expression in the name of the entry stands for the series made of these stable coefficients, and part (c) says that this series is
$A$. So the definition in the name of the entry and formula (E1) agree, and the sequence A392210 is $a(n)=\coeff{n}{A}$.

\begin{theorem}\label{thm:main}
For every $n>1$, the coefficient $a(2n-1)$ is even.
\end{theorem}

\section{Proofs}\label{sec:proofs}

\para{The idea in brief.} The right-hand side of \eqref{eq:def} contains $A$ only inside a square, and squaring
``pushes'' differences to higher degree; this gives uniqueness. Substituting $M_j$ into \eqref{eq:def} and using
$M_1(M_j)=M_{j+1}$ gives the nested form, and each level of nesting fixes one more coefficient. Finally, modulo $2$ the
square of a series has only even powers of $x$, so the only odd power of $x$ in $A$ is $x$ itself.

\begin{lemma}[An equation with a square]\label{lem:square}
Let $T$ and $P$ be normalized series with coefficients in a commutative ring $R$. There is exactly one normalized
series $G$ with coefficients in $R$ such that
\[
G(x)=T(x)+\bigl(G(P(x))\bigr)^2 .
\]
If $T$ and $P$ have integer coefficients, then so does $G$, and reducing $G$ modulo $m$ gives the solution of the same
equation over $\Z/m\Z$ (with $T$ and $P$ reduced modulo $m$).
\end{lemma}

\begin{proof}
The point is that the coefficient of $x^n$ on the right-hand side depends only on the coefficients of $G$ below degree
$n$. To see this, let $G$ and $H$ be two series with constant term $0$ that agree below degree $r$, so $G-H=O(x^r)$.
Since $P$ starts with $x$, the power $P^k$ has no terms below $x^k$, and so $G(P)-H(P)=O(x^r)$ as well. Moreover
$G(P)+H(P)$ has constant term $0$. Hence
\[
G(P)^2-H(P)^2=\bigl(G(P)-H(P)\bigr)\bigl(G(P)+H(P)\bigr)=O(x^{r+1}).
\]
In other words, changing $G$ in degree $r$ and above changes the right-hand side only in degree $r+1$ and above.

So the equation fixes the coefficients of $G$ one at a time. The constant term is $0$, the coefficient of $x$ is the
coefficient of $x$ in $T$, which is $1$, and for $n\ge2$ the coefficient of $x^n$ in $G$ must equal the coefficient of
$x^n$ in $T+G(P)^2$, which has already been computed from the earlier coefficients. Only additions and multiplications
are used, so this works in every commutative ring, and gives integer coefficients when $T$ and $P$ have them. Finally,
reducing the integer solution modulo $m$ gives a normalized series that satisfies the reduced equation (reduction is
compatible with all the operations involved), and by uniqueness it is \emph{the} solution over $\Z/m\Z$.
\end{proof}

\begin{proof}[Proof of Theorem~\ref{thm:def}]
(a) This is Lemma~\ref{lem:square} with $T=x$ and $P=M_1$.

(b) Since $M_0(x)=x$, we have $G_0=A$. Substitute $M_j(x)$ for $x$ in \eqref{eq:def}. Substitution can be done in steps,
and substituting into a square is the same as squaring the substituted series, so
\[
A(M_j)=M_j+\Bigl(A\bigl(M_1(M_j)\bigr)\Bigr)^2=M_j+A(M_{j+1})^2,
\]
using $M_1(M_j)=M_{j+1}$ from \eqref{eq:mobius}. This is $G_j=M_j+G_{j+1}^2$.

(c) Fix $d\ge1$, and set $H_{d+1}=0$, $H_j=H^{(d)}_j$ for $1\le j\le d$, and $H_0=A^{[d]}$. Then
$H_j=M_j+H_{j+1}^2$ for $0\le j\le d$, the same rule as in part (b). We show by downward induction on $j$ that
\[
G_j-H_j=O\bigl(x^{\,d+2-j}\bigr)\qquad(0\le j\le d+1).
\]
For $j=d+1$ this says $G_{d+1}=O(x)$, which is true because $G_{d+1}$ has constant term $0$. If the claim holds for
$j+1$, then
\[
G_j-H_j=G_{j+1}^2-H_{j+1}^2=\bigl(G_{j+1}-H_{j+1}\bigr)\bigl(G_{j+1}+H_{j+1}\bigr)=O\bigl(x^{\,d+1-j}\bigr)\cdot O(x),
\]
because $G_{j+1}$ and $H_{j+1}$ both have constant term $0$. This is $O(x^{d+2-j})$, as claimed. For $j=0$ we get
$A-A^{[d]}=O(x^{d+2})$.
\end{proof}

\begin{proof}[Proof of Theorem~\ref{thm:main}]
Reduce \eqref{eq:def} modulo $2$; this is allowed because reduction is compatible with all the operations involved. By
Lemma~\ref{lem:frobenius}, for every series $U$ with coefficients in $\F_2$ we have $U(x)^2=U(x^2)$, which contains only
even powers of $x$. Applying this to $U=A(M_1(x))$, we see that, modulo $2$, the right-hand side $x+A(M_1)^2$ has no odd
powers of $x$ except $x$ itself. Hence the coefficient of $x^{2n-1}$ in $A$ is even whenever $2n-1\ge3$, that is,
whenever $n>1$.
\end{proof}

\section{Checking the results}\label{sec:checks}

A mathematical proof can contain mistakes, so the results were checked in two independent ways.

\para{Numerical check.}
We computed $A$ by repeating the substitution $A\mapsto x+A(M_1)^2$ with exact integer arithmetic, starting
from $A=x$. This reproduces all 24 terms printed in the OEIS entry. We also evaluated the nested expression in the name of the entry,
cut off after $202$ levels, and found the same coefficients below degree $200$, as Theorem~\ref{thm:def}(c) predicts.
All coefficients of odd degree $2n-1$ with $3\le2n-1<200$ are even.
The computation was done with the PARI/GP computer algebra system, using code that was written separately from the
formal proofs (file \texttt{verification/independent\_checks.gp} in the accompanying repository~\cite{repo}).
A computation can only test finitely many cases, so it is a check on the definitions and on the indexing, not a proof.

\para{Formal proof in Lean.}
Lean~\cite{lean} is a \emph{proof assistant}: a program in which mathematical statements and their proofs are written
in a precise formal language, and every step is checked mechanically down to the basic rules of logic. The theorems of Section~\ref{sec:results} are proved in Lean~4 (version 4.19.0) on top of
Mathlib~\cite{mathlib}, Lean's library of formalized mathematics; the steps proved only on paper are identified below. The statements are collected in the file \texttt{lean/Entries/A392210.lean} of the repository,
in the namespace \texttt{OEIS.A392210}. There $x/(1-x)$ is called \texttt{geometric}, and the proofs are in the module \texttt{lean/Families/NestedBinomial.lean}. Parts (b) and (c) of Theorem~\ref{thm:def}, which connect formula (E1) with the nested expression in the name of the entry, are proved in this paper but are not part of the Lean formalization.
The table lists where each result of this paper is proved.
\begin{center}\small\renewcommand{\arraystretch}{1.25}
\begin{tabular}{>{\raggedright\arraybackslash}p{0.52\linewidth}>{\raggedright\arraybackslash}p{0.4\linewidth}}
\toprule
Statement in this paper & Lean theorem(s)\\
\midrule
Theorem~\ref{thm:def}(a) for integer coefficients & \texttt{third\_\allowbreak{}existence} \\
$A$ is normalized and satisfies \eqref{eq:def}; and $M_1(x)\cdot(1-x)=x$ & \texttt{third\_\allowbreak{}definition} \\
Theorem~\ref{thm:main}, for $A$ and for any normalized integer solution of \eqref{eq:def} & \texttt{third\_\allowbreak{}evenness}, \texttt{third\_\allowbreak{}arbitrary\_\allowbreak{}source} \\
\bottomrule
\end{tabular}
\end{center}
The Lean project was rebuilt from source, and the axioms on which each of these theorems depends were listed with
Lean's \texttt{\#print axioms} command. Every one depends only on the three standard axioms of Lean's logic
(\texttt{propext}, \texttt{Classical.choice}, \texttt{Quot.sound}); no step is left unproved. A formal proof certifies
exactly the statement that Lean checks. Whether that statement says what the OEIS entry means is a separate question,
so the Lean definitions were also compared with the definition in the OEIS entry
(Section~\ref{sec:sequence}).

\section{Discussion}\label{sec:discussion}

\para{Scope.} Theorem~\ref{thm:main} proves the observation of the entry for every $n$. It is an elementary fact, and the
point of this paper is to give a complete and checked account of it, including the meaning of the nested expression.
The theorem says nothing about the terms of even index, nor about remainders modulo $4$ or other numbers.

\para{The even-indexed terms.} The same idea gives a rule for the parity of the remaining terms. Apply
Lemma~\ref{lem:frobenius} to $U=A(M_1(x))$ more carefully: modulo $2$, $U(x)^2=U(x^2)=A\bigl(M_1(x^2)\bigr)$. So
\[
A(x)\equiv x+A\Bigl(\frac{x^2}{1-x^2}\Bigr)\pmod 2 .
\]
Comparing the coefficients of $x^{2m}$ on both sides gives, for every $m\ge1$,
\[
a(2m)\equiv\sum_{k=1}^{m}\binom{m-1}{k-1}\,a(k)\pmod 2 .
\]
So the parity of $a(2m)$ is determined by the parities of $a(1),\dots,a(m)$. (This rule is not part of the Lean
formalization; we checked it numerically for all $2m<130$.) Unlike the odd-indexed terms, the even-indexed terms show no
simple pattern in our data.

\para{Other equations of the same kind.} The entry lists further identities of the same shape, for example
$A\bigl(x/g(x)\bigr)=x/g(x)+A\bigl(x/(g(x)-x)\bigr)^2$ for every series $g$ with constant term
$1$. All of them follow from formula (E1) by substitution, in the same way as part (b) of Theorem~\ref{thm:def}.

\section*{About this paper}\label{sec:about}
\para{How the work was done.}
This paper is one of 29 papers on conjectures from the OEIS that came out of an AI research effort directed by
Dev~D.~Goyal~\cite{repo}. The roles were as follows.
\begin{itemize}
\item \emph{Dev~D.~Goyal} started the project, chose which problems to pursue, asked for proofs, formal verification and
revisions, and decided what to publish. Dev~D.~Goyal did not supply proof strategies or intermediate arguments.
\item \emph{ChatGPT} (OpenAI), using the GPT-6 model with Pro reasoning, worked autonomously from those prompts. It
searched the OEIS for suitable conjectures, found the proofs, wrote the Lean~4 formalization, and wrote a first draft
of this paper.
\item \emph{Claude Opus~5.5}, used with Claude Code, checked the work independently. It rebuilt the Lean project from source and confirmed that the theorems listed in Section~\ref{sec:checks} rely only on
Lean's standard axioms; compared the Lean definitions with the OEIS definition; recomputed the sequence with separately
written PARI/GP code, compared it with the terms published in the OEIS, and tested every claim of the paper on many
terms; and read every step of the proofs. It then wrote the present version of the paper, expanding the
explanations and, in places, writing out an argument in more detail or by a more elementary route. The theorems of
Section~\ref{sec:results} and the proof strategies are ChatGPT's. The Discussion section, including the short arguments
in it, was written by Claude Opus~5.5; its observations were checked by computer and are not part of the Lean
formalization.
\end{itemize}
Dev~D.~Goyal has verified the results and is responsible for this paper.

\para{Review status.} This paper has not been peer reviewed, and it has not been submitted to a journal or to the OEIS.

\para{Priority.} An OEIS entry that labels a statement as a conjecture shows only that its contributor did not know
a proof; it does not show that no proof existed. During the research only limited searches for earlier work were
made. No claim of priority is made.

\para{Corrections.} Please report any error in the issue tracker of the repository~\cite{repo}, where the Lean sources and the numerical checks are also
available.

\section*{Acknowledgements}
The sequence, its formulas and the observation are due to Paul~D.~Hanna~\cite{oeis}.

\begin{thebibliography}{99}\small
\raggedright
\bibitem{oeis} P.~D.~Hanna, entry A392210 (G.f.: $x+\mathrm{sq}(x/(1-x)+\mathrm{sq}(x/(1-2x)+\cdots))$, an infinitely nested square), \emph{The On-Line Encyclopedia of Integer Sequences}, February 2026. \url{https://oeis.org/A392210}. Text as of October 2026.
\bibitem{oeisdb} N.~J.~A.~Sloane and The OEIS Foundation Inc., \emph{The On-Line Encyclopedia of Integer Sequences}. \url{https://oeis.org}.
\bibitem{stanley} R.~P.~Stanley, \emph{Enumerative Combinatorics}, Volume~2, Cambridge University Press, 1999. (Chapter~5 treats formal power series, composition and compositional inverses.)
\bibitem{niven} I.~Niven, Formal power series, \emph{American Mathematical Monthly} 76 (1969), 871--889.
\bibitem{lean} L.~de~Moura and S.~Ullrich, The Lean~4 theorem prover and programming language, in \emph{Automated Deduction -- CADE 28}, Lecture Notes in Computer Science 12699, Springer, 2021, 625--635.
\bibitem{mathlib} The mathlib Community, The Lean mathematical library, in \emph{Proceedings of the 9th ACM SIGPLAN International Conference on Certified Programs and Proofs (CPP 2020)}, 367--381. Version used: commit \texttt{c44e0c8ee63c}.
\bibitem{repo} D.~D.~Goyal, \emph{Proofs and corrections of OEIS conjectures}, repository, 2026. \url{https://github.com/dvgyl/oeis-a396843-proof}.
\end{thebibliography}
\end{document}
